\documentclass[11pt]{article}

\usepackage[a4paper,margin=2.5cm]{geometry}
\usepackage{amsmath,amssymb,bm}
\usepackage{physics}

\title{Minimal Derivation of the RHF--MP2 Second-Order Energy}
\author{}
\date{}

\begin{document}
\maketitle

\section*{Setup: Hamiltonian and M{\o}ller--Plesset partition}

Let $\ket{\Phi_0}$ be a closed-shell RHF Slater determinant constructed from
canonical \emph{spatial orbitals} $\{\phi_p\}$ satisfying
\begin{equation}
\hat f\,\phi_p = \varepsilon_p\,\phi_p,
\end{equation}
where $\hat f$ is the RHF Fock operator and $\varepsilon_p$ are the canonical
orbital energies.

In second quantization, summing over spin, the electronic Hamiltonian is
\begin{equation}
\hat H
=
\sum_{pq} h_{pq}\, E_{pq}
+
\frac{1}{2}\sum_{pqrs} \langle pq | rs\rangle\,
\bigl(E_{pr}E_{qs} - \delta_{qr}E_{ps}\bigr),
\end{equation}
where $E_{pq} = \sum_\sigma a_{p\sigma}^\dagger a_{q\sigma}$ is the
spin-summed excitation operator, and
\begin{equation}
\langle pq | rs\rangle
=
\iint \phi_p^*(1)\phi_q^*(2)\,\frac{1}{r_{12}}\,
\phi_r(1)\phi_s(2)\,d1\,d2
\end{equation}
are the (non-antisymmetrized) two-electron integrals in physicist notation.

We employ the M{\o}ller--Plesset partition
\begin{equation}
\hat H = \hat H_0 + \hat V,
\qquad
\hat H_0 = \sum_p \varepsilon_p\, E_{pp} \; (+\text{constant}),
\end{equation}
where the additive constant is chosen such that
\begin{equation}
E^{(0)} = \mel{\Phi_0}{\hat H_0}{\Phi_0} = E_{\mathrm{HF}}.
\end{equation}

\subsection*{Extension to electron--positron systems}

For a closed-shell electron reference plus a single positron, we use
barred indices $\bar p, \bar q, \ldots$ for positron orbitals and
ordinary Latin indices $p, q, \ldots$ for electrons.

The reference $\ket{\Phi_0}$ is constructed from canonical orbitals
satisfying the coupled Fock equations
\begin{equation}
\hat f^{(e)}\,\phi_p = \varepsilon_p\,\phi_p,
\qquad
\hat f^{(p)}\,\phi_{\bar p} = \varepsilon_{\bar p}\,\phi_{\bar p}.
\end{equation}

The Hamiltonian is
\begin{equation}
\hat H
=
\sum_{pq} h_{pq}\, E_{pq}
+
\frac{1}{2}\sum_{pqrs} \langle pq | rs\rangle\,
\bigl(E_{pr}E_{qs} - \delta_{qr}E_{ps}\bigr)
+
\sum_{\bar p\bar q} h_{\bar p\bar q}^{(p)}\, a_{\bar p}^\dagger a_{\bar q}
-
\sum_{pq\bar r\bar s} \langle p\bar r | q\bar s \rangle\,
E_{pq}\, a_{\bar r}^\dagger a_{\bar s},
\end{equation}
where $h^{(p)}$ differs from $h$ by the sign of the nuclear attraction,
and the electron--positron interaction is attractive (minus sign).
Since there is only one positron, there is no positron--positron two-body term.

The M{\o}ller--Plesset partition is
\begin{equation}
\hat H = \hat H_0 + \hat V,
\qquad
\hat H_0 = \sum_p \varepsilon_p\, E_{pp}
+ \sum_{\bar p} \varepsilon_{\bar p}\, a_{\bar p}^\dagger a_{\bar p}
\; (+\text{constant}),
\end{equation}
with $E^{(0)} = E_{\mathrm{HF}}$.

\section*{Zeroth-order eigenstates and denominators}

\subsection*{Electron double excitations}

For RHF with spatial orbitals, we consider spin-adapted double excitations.
Define the spin-summed doubles amplitudes $T_{ij}^{ab}$ which correspond to
excitations from doubly-occupied spatial orbitals $i,j$ to virtual spatial
orbitals $a,b$.
Since $\hat H_0$ is diagonal in the canonical basis, the eigenvalue shift is
\begin{equation}
\varepsilon_a + \varepsilon_b - \varepsilon_i - \varepsilon_j.
\end{equation}
We define the MP2 energy denominator
\begin{equation}
\Delta_{ij}^{ab} \equiv \varepsilon_i + \varepsilon_j - \varepsilon_a - \varepsilon_b.
\end{equation}

\subsection*{Extension to electron--positron excitations}

For electron--positron excitations, one electron from spatial orbital $i$
is excited to virtual orbital $a$, and the positron from orbital $\bar\imath$
is excited to virtual positron orbital $\bar a$.
Since $\hat H_0$ is diagonal in the canonical basis, the eigenvalue shift is
\begin{equation}
\varepsilon_a + \varepsilon_{\bar a} - \varepsilon_i - \varepsilon_{\bar\imath}.
\end{equation}
We define the electron--positron MP2 energy denominator
\begin{equation}
\Delta_{i\bar\imath}^{a\bar a} \equiv \varepsilon_i + \varepsilon_{\bar\imath} - \varepsilon_a - \varepsilon_{\bar a}.
\end{equation}

\section*{First-order wave function and doubles amplitudes}

\subsection*{Electron doubles}

Using intermediate normalization $\braket{\Phi_0}{\Psi}=1$, expand
\begin{equation}
\ket{\Psi} = \ket{\Phi_0} + \ket{\Psi^{(1)}} + \cdots,
\qquad
E = E^{(0)} + E^{(1)} + E^{(2)} + \cdots .
\end{equation}
The first-order correction satisfies, because $E^{(1)} = \mel{\Phi_0}{\hat V}{\Phi_0} = 0$,
\begin{equation}
(\hat H_0 - E^{(0)})\ket{\Psi^{(1)}} = -\,\hat V\ket{\Phi_0},
\end{equation}
projected onto the space orthogonal to $\ket{\Phi_0}$.

For canonical HF orbitals, Brillouin's theorem implies that singles do not
couple to the reference, so $\ket{\Psi^{(1)}}$ contains only double excitations.
The first-order doubles amplitudes are
\begin{equation}
T_{ij}^{ab}
=
\frac{\langle ij | ab\rangle}{\varepsilon_i + \varepsilon_j - \varepsilon_a - \varepsilon_b},
\end{equation}
where $\langle ij | ab\rangle$ is the non-antisymmetrized two-electron integral
over spatial orbitals.

\subsection*{Extension to electron--positron doubles}

Brillouin's theorem also applies to electron and positron singles separately.
The first-order wave function now includes electron--positron excitations
with amplitudes
\begin{equation}
T_{i\bar\imath}^{a\bar a}
=
\frac{-\langle i\bar\imath | a\bar a\rangle}{\varepsilon_i + \varepsilon_{\bar\imath} - \varepsilon_a - \varepsilon_{\bar a}},
\end{equation}
where the minus sign arises from the attractive electron--positron interaction.

\section*{Second-order energy}

\subsection*{Electron contribution}

The second-order correction to the energy is
\begin{equation}
E^{(2)} = \mel{\Phi_0}{\hat V}{\Psi^{(1)}}.
\end{equation}
For closed-shell RHF, summing over spin gives the well-known expression.
Substituting the MP2 amplitudes yields the canonical RHF--MP2 expression
\begin{equation}
\boxed{
E^{(2)}_{\mathrm{RHF}}
=
\sum_{ijab}
\frac{\langle ij | ab\rangle\bigl(2\langle ij | ab\rangle - \langle ij | ba\rangle\bigr)}
{\varepsilon_i + \varepsilon_j - \varepsilon_a - \varepsilon_b}
}.
\end{equation}
Here the factor of 2 arises from the spin summation (same-spin and opposite-spin
contributions), and the exchange term $\langle ij | ba\rangle$ accounts for
Fermi correlation among same-spin electrons.

\subsection*{Extension to electron--positron contribution}

The second-order correction now includes both contributions:
\begin{equation}
E^{(2)} = E^{(2)}_{ee} + E^{(2)}_{ep}.
\end{equation}
For the electron--positron term, there is no exchange between the distinguishable
particles, but there is a factor of 2 from summing over the two electron spin
orientations. Substituting the electron--positron MP2 amplitudes yields
\begin{equation}
\boxed{
E^{(2)}_{ep}
=
2\sum_{i\bar\imath a\bar a}
\frac{\big|\langle i\bar\imath | a\bar a\rangle\big|^2}
{\varepsilon_i + \varepsilon_{\bar\imath} - \varepsilon_a - \varepsilon_{\bar a}}
}.
\end{equation}

\end{document}
